\section{Design}\label{sec:design}

\subsection{Profile, not protocol}

cA2A is a trust profile layered on the Agent2Agent (A2A) protocol \citep{a2a2026}, not a
competing transport. A2A moves tasks between agents and, through the Signed Agent Card,
answers exactly one question: did the domain owner issue this card. It does not establish
that a peer runs attested code, that a delegating agent actually holds the authority it
passes on, that a task payload is confidential to the intended peer, or that there is an
unbroken record of who delegated what to whom. cA2A adds that trust envelope around a
delegated task without redefining how the task is carried. It binds to A2A the way TRACE
\citep{trace2025} binds to the IETF RATS \citep{rats2021}, EAT \citep{eat2023}, and SCITT
\citep{scitt2024} work: as an overlay, not a replacement. The profile is written to be
independent of any one organization, cloud, or TEE vendor; the implementation currently
appraises AMD SEV-SNP, Intel TDX, and TPM~2.0 evidence.

Four design tenets constrain every primitive below. First, \emph{profile, not protocol}:
reuse A2A and the RATS, EAT, and SCITT stack rather than reinvent them. Second, \emph{neutral by
construction}: no coupling to a single silicon vendor, cloud, or model platform. Third,
\emph{fail closed}: absence of evidence is denial, not a warning; a missing attestation or an
unverifiable chain denies the call. Fourth, \emph{offline verifiable}: a third party can
verify a delegation chain and its provenance DAG without contacting or trusting any operator.

\subsection{Where cA2A attaches to A2A}

cA2A interposes on the trust decision that A2A leaves undefined: the moment one agent
receives a delegated task from a peer and must decide whether to accept it. On an inbound
peer call, the caller appraises the callee's attestation offer and seals the payload to the
channel key it vouches for. The callee then, in order, (i) verifies the delegation chain
(trusted root, signature, validity window, continuity, attenuation, in-chain replay, and
revocation when a snapshot is supplied), (ii) checks a holder proof showing that the caller
holds the key the chain delegates to, (iii) intersects the delegated scope with local policy,
(iv) appraises the caller's own attestation when the callee requires it, (v) decides the
requested capability, and (vi) only then opens the sealed payload and emits a linked
provenance record for the allow or deny decision. The A2A transport, agent identity issuance beyond delegation, and model governance
are all out of scope: cA2A governs only the delegation path between agents.

\subsection{Threat model}

cA2A defends the delegation path between agents against a capable adversary who may operate a
peer agent in a trust domain cA2A does not control; present a valid A2A Signed Agent Card
while running tampered or unmeasured code; sit on the network between two agents or operate
the host a peer runs on; and attempt to widen a delegated grant, replay a credential into
another chain, or reparent a provenance record. The assets and their defenses are summarized
in Table~\ref{tab:primitives}.

Out of adversary scope is breaking the underlying cryptographic primitives (Ed25519
\citep{eddsa2017}, the hash function) or compromising TEE firmware or hardware microcode
\citep{amdsevsnp2023,inteltdx2023}. Those are the trust anchors. The central attack cA2A
addresses is the confused deputy \citep{confuseddeputy1988} at machine speed: a delegate that
accepts a narrow task and then acts with authority it was never granted.

What ca2a 0.3.1 defends is narrower than this threat model. Attenuated delegation,
scope-policy intersection on the inbound path, the sealed channel's cryptography, and signed
provenance ship and are tested in software (Section~\ref{sec:evaluation}). Attestation
verifiers for SEV-SNP, TDX, and TPM evidence ship, and the live path can require hardware
assurance, but this is opt in: the reference client and server default to software assurance.
Hardware runs so far are one-directional across two clouds or mutual within one operator
(Section~\ref{subsec:attestation}). Confidentiality against a hosting operator additionally
depends on platform properties the recorded hosts lacked (Section~\ref{sec:after}). cA2A
should therefore not be described as attested across independent trust domains, and task
payloads should not be assumed protected from a hosting operator.

\subsection{Primitive 1: attenuated delegation}\label{subsec:delegation}

A delegation chain is a root-to-leaf list of signed credentials that lets Agent~A hand a
bounded slice of its authority to~B, and B a still-smaller slice to~C, with each grant
provably within the one above it. The semantics mirror the signed delegation chain
implemented and tested in agent-manifest \citep{agentmanifest2026}, and follow the
capability-attenuation model of macaroons \citep{macaroons2014}; cA2A implements that verifier
and its runtime calls it on every inbound peer request.

Each hop carries a \texttt{DelegationCredential} with a signed body and a detached Ed25519
signature. The body has \texttt{credential\_id} (a unique id for the hop), \texttt{issuer}
and \texttt{subject} (the delegator's and delegate's public keys), \texttt{scope} (the set of
capabilities granted at this hop), \texttt{depth} (0 at the root, $+1$ per hop),
and \texttt{parent\_id} (the parent hop's \texttt{credential\_id}, null at the root), with
optional \texttt{not\_before} and \texttt{not\_after} bounds. The signed bytes are the RFC~8785
\citep{jcs2020} canonical form of the body, with \texttt{scope} as a sorted array; floats are
refused and integers are bounded to the range RFC~8785 can represent exactly.

The verifier, \texttt{verify\_chain}, raises the specific error for the first invariant that
fails. The root issuer must be in the verifier's trusted set (else
\texttt{UNTRUSTED\_DELEGATION\_ROOT}); every hop's signature must verify against its issuer
(else \texttt{INVALID\_CREDENTIAL}) and its validity window must contain the evaluation time;
the root must have no parent and depth 0, each hop's \texttt{parent\_id} must equal the
previous hop's \texttt{credential\_id}, and each hop's \texttt{issuer} must equal the previous
hop's \texttt{subject}, and each hop's depth must be the previous depth $+1$ (else
\texttt{BROKEN\_DELEGATION\_LINK}); no depth may exceed \texttt{max\_depth} (else
\texttt{DELEGATION\_DEPTH\_EXCEEDED});
each hop's scope must be a subset of its parent's scope (else \texttt{SCOPE\_ESCALATION});
and no \texttt{credential\_id} may repeat (else \texttt{CREDENTIAL\_REPLAY}). Given a signed
revocation snapshot, a hop revoked by its issuer or any issuer above it is refused
(\texttt{CREDENTIAL\_REVOKED}). The verifier keeps no global history of used credentials, so
replay of a whole valid chain by a different party is addressed by the holder proof in the
inbound path rather than by \texttt{verify\_chain}.

Consider the chain
\[
  \text{A} \xrightarrow{\ S_A\ } \text{B} \xrightarrow{\ S_B \subseteq S_A\ } \text{C},
  \qquad S_C \subseteq S_B \subseteq S_A .
\]
A issues B a child credential whose scope $S_B$ is a subset of A's grant $S_A$, signed over
the canonical body. B in turn issues C a credential with $S_C \subseteq S_B$. The subset
check at every hop is the confused-deputy defense within a chain: it forecloses B accepting a narrow task
from A and then acting with authority A never granted. Attenuation is the whole point.
The \texttt{parent\_id}, \texttt{issuer}/\texttt{subject} continuity, and unique-id checks
close the complementary attack of replaying a credential within a chain or splicing one from
another chain into it. They do not bind a chain to a particular task or execution: a credential
carries no task or execution identifier, so the delegate can present the same valid chain
with a different request. Making execution lineage part of the authorization decision is the
subject of history-based access control \citep{abadi2003history}, provenance-based access control
\citep{park2012pbac}, and, for agent delegation, the PIC model \citep{gallo2026pic}. cA2A does not
provide it; authenticated lineage records (Section~\ref{sec:after}) say which hop acted under
which credential, which is not the same as refusing authority outside its lineage. This primitive is implemented and validated (Properties~\ref{prop:c1}
and~\ref{prop:c2}).

\subsection{Primitive 2: runtime peer attestation}\label{subsec:attestation}

Before a peer is trusted with a delegated task, it presents evidence of the code it runs. cA2A
follows the TEE provider pattern of cMCP \citep{cmcp2026}. A provider
implements \texttt{detect()}, returning whether it is available on the current host, and
\texttt{attest(public\_key, nonce)}, returning an \texttt{AttestationReport} that binds a
public key to the host's hardware measurement under a nonce. The report carries
\texttt{platform}, \texttt{measurement}, the bound \texttt{public\_key}, and the \texttt{nonce},
following the RATS evidence model \citep{rats2021} and EAT claims format \citep{eat2023}.

The providers are a \texttt{software-only} backend (for development and CI, which reports
\texttt{platform:\ software-only} and never a hardware platform string) and three hardware
backends: \texttt{tpm} (TPM 2.0 or vTPM), \texttt{sev-snp} (AMD SEV-SNP \citep{amdsevsnp2023}),
and \texttt{tdx} (Intel TDX \citep{inteltdx2023}).

\textbf{Implementation status at 0.3.1.} The SEV-SNP and TDX verifiers check certificate chains
to pinned AMD and Intel roots, report signatures, and key binding, and the project records runs
against genuine evidence from both platforms; for this edition the TDX verifier was re-run
against a public GCP TDX quote (Section~\ref{sec:evaluation}). The TPM verifier checks quote
structure, binding, and attestation-key signature; its certificate chain depends on a
caller-supplied vendor root, and on one Azure host generation no chain to a pinnable root could
be assembled. SEV-SNP and TDX collectors use the kernel configfs-TSM interface and have run on
GCP confidential VMs; an Azure SEV-SNP guest is paravisor-mediated, cannot set
\texttt{REPORT\_DATA} directly, and roots its channel key through the vTPM instead. The project
records two live runs. On July 27, 2026 an Azure SEV-SNP caller appraised a GCP TDX peer's quote
against the Intel root, sealed a delegated task to the bound key, and the peer allowed an
in-scope capability and refused an out-of-scope one; the callee did not appraise the caller, and
one operator ran both sides. On September 17, 2026 two GCP SEV-SNP guests in one project
appraised each other and exchanged sealed non-secret payloads in both directions. Neither run is
mutual attestation between independent operators, so Property~\ref{prop:c6} is not validated on
hardware. Verification fails closed when required evidence is absent. In 0.3.1 a required caller
appraisal could be satisfied by an offer relayed from another node (Section~\ref{sec:after}).

\subsection{Primitive 3: sealed peer channel}\label{subsec:sealed}

When A sends B a task payload, that payload crosses a network and lands in B's memory. If B is
in another trust domain, mTLS protects the pipe but not the endpoint: the operator hosting B,
or a connectivity provider between them, can read plaintext. The sealed channel encrypts the
payload to B's channel public key, and B's attestation evidence binds that key to a
measurement, extending the attestation-gated key pattern cMCP uses for the agent-to-gateway
case to the peer-to-peer case. The connectivity path sees ciphertext. The interface is
\texttt{SealedChannel(peer\_public\_key).seal} on the sender and \texttt{open\_sealed} with the
private key on the receiver.

\textbf{Cryptographic layer: implemented and validated.} The channel ships over
X25519 ECDH, HKDF-SHA256 key derivation, and ChaCha20-Poly1305 AEAD. A payload sealed to a
peer's channel key decrypts only with that key, the connectivity path sees only ciphertext, and
a tampered blob fails AEAD authentication so \texttt{open} fails closed rather than returning
plaintext. Property~\ref{prop:c4} validates exactly this at the cryptographic layer
(Section~\ref{sec:evaluation}).
\textbf{Binding to attestation.} On the live path the caller seals only to a channel key that
its verifier has appraised (\texttt{ca2a\_runtime.attestation}). In software mode that appraisal
has assurance \texttt{none}; with a hardware verifier configured and
\texttt{require\_hardware} set, a software offer is refused before sealing. Even with hardware
appraisal, the channel key is an ordinary key held by a process in the guest. In the
September 17 diagnostic the launch measurements did not cover application code installed after
boot, and the SEV-SNP hosts reported that ciphertext hiding was not enabled. Confidentiality
against the hosting operator therefore remains the dependency that Property~\ref{prop:c4}
names.

Figure~\ref{fig:sealed-channel} shows the seal-and-open path. The sender derives a
one-time symmetric key from an ephemeral X25519 ECDH to the peer's channel public key,
runs it through HKDF-SHA256, and encrypts the payload under ChaCha20-Poly1305. The
connectivity path sees only the sealed blob; only the holder of the peer's channel private key
recovers the plaintext.

\begin{figure}[h]
\centering
\begin{tikzpicture}[
  font=\small,
  node distance=6mm,
  cryptobox/.style={draw, rounded corners, minimum height=8mm, align=center,
    fill=blue!6, inner sep=3pt},
  enclavebox/.style={draw, rounded corners, minimum height=8mm, align=center,
    fill=green!8, inner sep=3pt},
  flow/.style={-{Stealth[length=2mm]}, thick},
]
  % Sender-side derivation pipeline.
  \node[cryptobox] (eph) {ephemeral\\X25519};
  \node[cryptobox, right=of eph] (ecdh) {ECDH\\to peer key};
  \node[cryptobox, right=of ecdh] (hkdf) {HKDF-\\SHA256};
  \node[cryptobox, right=of hkdf] (aead) {ChaCha20-\\Poly1305};
  \draw[flow] (eph) -- (ecdh);
  \draw[flow] (ecdh) -- (hkdf);
  \draw[flow] (hkdf) -- (aead);

  % Untrusted path carrying ciphertext.
  \node[draw, dashed, rounded corners, below=10mm of hkdf, minimum width=32mm,
    minimum height=7mm, align=center, fill=red!5] (wire)
    {untrusted path / host\\(sees ciphertext only)};
  \draw[flow] (aead.south) |- (wire.east) node[midway, above right, font=\scriptsize]
    {sealed blob};

  % Peer enclave that opens the blob.
  \node[enclavebox, below=10mm of eph, minimum width=26mm] (peer)
    {peer TEE guest\\open with\\private key};
  \draw[flow] (wire.west) -- (wire.west -| peer.east);
  \node[below=1mm of peer, font=\scriptsize, align=center]
    {key bound to appraised\\evidence (hardware opt-in)};

  % Grouping label for the sender pipeline.
  \begin{scope}[on background layer]
    \node[draw=blue!30, rounded corners, fit=(eph)(aead), inner sep=5mm,
      label=above:{sender: seal to peer public key}] {};
  \end{scope}
\end{tikzpicture}
\caption{Sealed peer channel. An ephemeral X25519 ECDH to the peer's channel key is run
through HKDF-SHA256 to derive a symmetric key, and the payload is sealed under
ChaCha20-Poly1305 AEAD. The untrusted connectivity path sees only ciphertext; only the
private key matching the appraised channel key opens the blob. Whether the peer's host can
reach that key depends on the platform (Section~\ref{sec:after}).}
\label{fig:sealed-channel}
\end{figure}

\subsection{Primitive 4: provenance DAG}\label{subsec:provenance}

cA2A emits a signed TRACE record per delegation hop \citep{trace2025}. The A2A profile adds a
delegation-link block to the base TRACE record, which otherwise carries a subject but no notion
of a parent. The block has \texttt{delegation.parent\_record\_hash}, the hash of the parent
hop's TRACE record, and \texttt{delegation.credential\_id}, the id of the delegation credential
this hop acted under. A root hop omits the block; each subsequent hop references its parent, so
the set of records across A~to~B~to~C reconstructs the delegation tree. The construction follows
the W3C PROV data model \citep{prov2013} for lineage and the SCITT \citep{scitt2024} approach to
transparency over signed statements.

Given the records for a workflow and the recipient's trusted signing keys, a verifier
(\texttt{verify\_trace\_dag}) checks each record's signature and rebuilds the parent links from
\texttt{parent\_record\_hash}; a cross-check (\texttt{cross\_check\_trace\_dag}) then aligns the
path with a separately verified delegation chain, requiring matching path length and credential
identifiers and that each record was signed by the key the chain names as that hop's delegate.
Tampering any signed record breaks its signature, and altering a parent link breaks the chain of
hashes. The runtime also keeps an unsigned \texttt{DelegationRecord} form whose hash links give
only structural consistency: a party able to rewrite the whole path can recompute every link,
so since ca2a 0.3.0 the CLI reports such a path as structurally consistent but not
authenticated. Property~\ref{prop:c5} and Section~\ref{sec:eval-c5} distinguish the two. Both
verifiers accept one ordered root-to-leaf path. The A2A profile's delegation-link block is
defined in trace-spec.

Figure~\ref{fig:delegation-dag} shows both halves together: the attenuated delegation chain
on top, where each hop's scope is a subset of its parent's (strictly smaller in this example), and the hash-linked
provenance path below, where each record carries the hash of its parent record and names the
credential the hop acted under. The two are bound: record $i$ must match credential $i$ and, for
signed records, must be signed by that credential's delegate key.

\begin{figure}[h]
\centering
\begin{tikzpicture}[
  font=\small,
  credbox/.style={draw, rounded corners, minimum width=30mm, minimum height=11mm,
    align=center, fill=blue!6, inner sep=2pt},
  recbox/.style={draw, rounded corners, minimum width=30mm, minimum height=10mm,
    align=center, fill=orange!8, inner sep=2pt},
  deleg/.style={-{Stealth[length=2mm]}, thick},
  hashlink/.style={-{Stealth[length=2mm]}, thick, densely dashed},
  bind/.style={dotted, thick, gray},
]
  % Attenuated delegation chain (top row): scope narrows per hop.
  \node[credbox] (a) {A $\to$ B\\scope $\{r,w,x\}$};
  \node[credbox, right=14mm of a] (b) {B $\to$ C\\scope $\{r,w\}$};
  \node[credbox, right=14mm of b] (c) {C $\to$ D\\scope $\{r\}$};
  \draw[deleg] (a) -- node[above, font=\scriptsize] {$\subseteq$} (b);
  \draw[deleg] (b) -- node[above, font=\scriptsize] {$\subseteq$} (c);

  % Provenance records (bottom row), hash-linked to parent.
  \node[recbox, below=16mm of a] (r0) {record 0\\parent: none};
  \node[recbox, below=16mm of b] (r1) {record 1\\parent hash $h_0$};
  \node[recbox, below=16mm of c] (r2) {record 2\\parent hash $h_1$};
  \draw[hashlink] (r1) -- node[below, font=\scriptsize] {$h_0$} (r0);
  \draw[hashlink] (r2) -- node[below, font=\scriptsize] {$h_1$} (r1);

  % Bind record i to credential i.
  \draw[bind] (a) -- (r0);
  \draw[bind] (b) -- (r1);
  \draw[bind] (c) -- (r2);

  \node[right=3mm of c, align=left, font=\scriptsize] {narrowing\\scope};
  \node[left=2mm of r0, align=right, font=\scriptsize] {hash-\\linked\\DAG};
\end{tikzpicture}
\caption{Attenuated delegation chain (top, solid) and its hash-linked provenance DAG
(bottom, dashed). Each delegation hop's scope is a provable subset ($\subseteq$) of its
parent's; each provenance record stores the SHA-256 hash of its parent record. The dotted
lines bind record $i$ to credential $i$. With signed records and trusted keys, tampering,
reparenting, or attributing a record to the wrong delegate is detected offline; hash links
alone give only structural consistency.}
\label{fig:delegation-dag}
\end{figure}

\subsection{Scope-policy intersection}\label{subsec:cedar}

Verifying that a delegated grant is well-formed is not the same as deciding that the local peer
is willing to act on it. On an inbound call, the cA2A runtime is specified to intersect the
delegated scope $S_B$ with the peer's local authorization policy, expressed in Cedar
\citep{cedar2023} or as a simple allow set, and to act only on the intersection
$S_B \cap P_{\text{local}}$. A grant is
thus bounded from above by both the delegator's authority (attenuation) and the delegate's own
policy, so neither a broad grant nor a permissive local default alone can widen what B actually
does.

\textbf{Implementation status.} The intersection runs against either a Cedar engine
(\texttt{cedarpy} 4.8.7) or the set-based \texttt{LocalPolicy}, inside the inbound decision
path (\texttt{handle\_peer\_request}). That path is served by the reference HTTP transport
(since ca2a commit \texttt{7d33445}, July 2026) and by a bridge to the official \texttt{a2a-sdk}
(commit \texttt{f26e2bc}), which the project tests over a local JSON-RPC loopback with
\texttt{a2a-sdk} 1.1.4. Property~\ref{prop:c3} is evaluated in Section~\ref{sec:eval-c3}.

\subsection{Primitives to properties}

Table~\ref{tab:primitives} maps each primitive to the asset it protects, the security property
it supports, and its status at ca2a 0.3.1. The software results need no hardware TEE and can be
rerun from the open-source codebase; the hardware dependencies are stated per row.

\begin{table}[h]
\centering
\caption{cA2A primitives, the assets they protect, the properties they support, and status.
``Validated'' means backed by the rerun software experiments of Section~\ref{sec:evaluation},
at the layer stated.}
\label{tab:primitives}
\small
\begin{tabular}{@{}>{\raggedright\arraybackslash}p{3.0cm}>{\raggedright\arraybackslash}p{3.4cm}>{\raggedright\arraybackslash}p{3.4cm}>{\raggedright\arraybackslash}p{3.4cm}@{}}
\toprule
\textbf{Primitive} & \textbf{Asset protected} & \textbf{Property} & \textbf{Status} \\
\midrule
Attenuated delegation & Bounded authority (no delegate exceeds its grant) &
  C1 attenuation soundness; C2 in-chain replay and splice rejection & Validated (software) \\
Scope-policy intersection & Bounded authority under local policy &
  C3 scope-policy intersection & Validated (software), including the reference inbound path \\
Runtime peer attestation & Peer integrity (only appraised code trusted) &
  C6 cross-operator attestation & Synthetic evidence only; hardware runs one-directional or same-operator \\
Sealed peer channel & Task confidentiality to the intended peer &
  C4 sealed-payload confidentiality & Validated (cryptographic layer); key confinement not shown \\
Provenance record & Provenance (signed record of who delegated what) &
  C5 provenance integrity & Structural checks validated; signed path covered by project tests \\
\bottomrule
\end{tabular}
\end{table}
